%!TEX root = ../../thesis.tex

\section{Fuzzing}\label{sec:fuzzing}

% Drafting outline: bullet points to be fleshed out into prose.
% \todo{cite: ...} marks claims whose canonical reference is not yet in references/*.bib.

\subsection{Overview and Motivation}\label{sec:fuzzing-overview}

% --- Verbatim excerpt from the proposal (own text; citation keys remapped to references/*.bib) ---
Fuzz testing, or \textit{fuzzing} for short, is a widely deployed method for automatic
software testing and vulnerability discovery that makes use of a diverse set of
tools~\cite{manes_art_2019}.
Its goal is to find inputs to a Program Under Test (PUT) that trigger a certain
behavior, often a crash that might then lead to the discovery of bugs or security
vulnerabilities.
% --- End excerpt ---

The technique was first introduced in the early 1990s as a way of testing the
reliability of UNIX utilities~\cite{miller_empirical_1990}.
The authors were inspired by spurious characters being introduced to a dial-up
line causing crashes of common programs on a remote workstation.
This lad to the creation of a simple command-line utility \emph{fuzz} that produced
random input strings that were then sent to various utilities in the hope of causing
a crash in the form of a core dump or a hang.
They managed to find inputs that caused between 24\% and 33\% of the tested
utility programs to crash or hang.
By analyzing these inputs, bugs could be identified and fixed, thereby
contributing to the reliability of the system.

The same approach is still in use today to find security vulnerabilities in
input-processing code, especially memory-safety bugs in C/C++.
For example, Google's OSS-Fuzz
\enquote{aims to make common open source software more secure and stable by combining
modern fuzzing techniques with scalable, distributed execution}~\cite{arya_oss-fuzz_nodate}.
It supports a variety of target languages and is therefore widely applicable.
OSS-Fuzz has helped to fix over 13,000 security vulnerabilities and 50,000 bugs across
1,000 projects, proving the importance of modern fuzzing techniques for ensuring reliability and security.\\

\subsection{Core Terminology\label{sec:fuzzing-definitions}}

We will introduce some general definitions that mainly follow \textcite{manes_art_2019}.

\begin{enumerate}
  \item \emph{Program Under Test (PUT):}
    The program under test is an executable software artifact like a command-line utility
    or library together with an interface that consumes input. In our case it denotes
    a binary executed through the MOGI emulator\cref{chap:mogi-setup}.

  \item \emph{Security Policy:}
    A security policy defines which observable behaviors of the PUT are acceptable.
    Each execution either satisfies or violates a given security policy. A common
    security policy deems all executions that result in crashes like segmentation
    fauls to be violations.

  \item \emph{Bug Oracle:}
    A bug oracle is a program that decides whether a given input to the PUT violates
    a given security policy. A widely used bug oracle might deem an execution
    a violation if the PUT crashes or hangs.

  \item \emph{Fuzzing:}
    Fuzzing denotes the process of feeding the PUT with inputs sampled from an input space.
    This input space may include expected and valid inputs as well as invalid, unexpected
    inputs in the hope of violating a given security policy.

  \item \emph{Fuzz Testing:}
    Fuzz testing denotes the software testing technique that makes use of fuzzing.

  \item \emph{Fuzzer:}
    A fuzzer is a program that performs fuzz testing on a PUT.

  \item \emph{Fuzz Campaign:}
    A fuzz campaign is a specific execution of a fuzzer on a PUT with a specific
    security policy. Its goal is to find bugs in the PUT that violate this policy.

  \item \emph{Test Case:}
    A test case denotes a concrete input that the fuzzer feeds to the PUT during
    a fuzzing campaign. The resulting execution is then judged by the bug oracle
    and the test case recorded in case the bug oracle deems it a violation of the
    security policy.

  \item \emph{Fuzz Configuration:}
    A fuzz configuration comprises the parameter values that control the fuzzing
    algorithm.
\end{enumerate}

% TODO include the Fuzz Testing ALGORITHM scheme with the steps PREPROCESS, CONTINUE, etc.

\subsection{Taxonomy\label{sec:fuzzing-taxonomy}}

Fuzzers can be roughly divided into three classes, defined by the amount of
information from the execution of the PUT available to them.~\cite{liang_fuzzing_2018}.

\subsubsection{Black-box fuzzing}

Black-box fuzzers have no access to the internal behavior of a PUT during execution.
They can only make use of the bug oracle to decide whether a given input violates
the security policy or not.
The core advantage of this approach is the wide applicability since no instrumentation
is required.
This lack of potential instrumentation overhead also benefits execution throughput.
The obvious disadvantage is the lack of information about PUT behavior.
This makes it difficult for fuzzers to generate high-coverage inputs.

\subsubsection{White-box fuzzing}

In constrast to pure black-box fuzzers, white-box fuzzers have full access to
the internal logic of the PUT as well as usually any original source code or
documentation which can greatly aid the fuzzing process.
In practice, this allows an approach known as \emph{symbolic execution}: instead
of relying on actual execution traces, a white-box fuzzer can gather symbolic
constraints at branching points and use constraint solving to find inputs that exercise
a specific path.
While this allows to easily find test cases of interest in theory, this approach scales poorly
to large PUTs due to the path explosion problem.
A popular white-box fuzzer is SAGE~\cite{TODO cite godefroid12}.

\subsubsection{Grey-box fuzzing}

Since full information is often not available but observing certain parts
of the internal behavior of the PUT might still be possible, \textit{grey-box}
fuzzing represents an important middle-ground between the previous two extremes.
By making use of runtime instrumentation information like code coverage, the fuzzer can
better judge the behavior of the PUT for a given input and direct its exploration efforts towards promising regions.
Limitations are discussed in detail in~\cref{sec:fuzzing-limitations}.

\subsubsection{Hybrid approaches}

There also exists a variety of hybrid approaches combining different fuzzing methodologies.
A prominent example is \emph{Driller} that alternates pure greybox fuzzing with
\emph{concolic} execution (a combination of concrete and symbolic execution) to overcome hard branch
conditions like CNC checks that have historically presented a major obstacle for fuzzers~\cite{stephens_driller_2016}

\subsection{Input generation}\label{sec:fuzzing-input-generation}

% Outline follows Manès et al. §5 (Input Generation); the § pointers below refer to that paper.

The content of a test case directly controls whether a bug is triggered or not, making
it one of the most influential design decisions of a fuzzer~\cite{manes_art_2019}.
Input generation can be categorized into either generation- or mutation-based.

\subsubsection{Generation-based input generation}

Generation-based input generation makes use of a model of the input expected by
the PUT, often represented through a grammar specification.
These models can either be \emph{predefined}, \emph{inferred}, or \emph{encoder-based}.

Predefined models are specified by the user.
They are often used by network fuzzers consuming protocol specifications or kernel fuzzers
consuming system call interfaces.

Inferred models don't rely on an explicitly specified model but rather infer it through some automatic means.
Inference can either take place as a preprocessing step or during a fuzz campaign.
Test-Miner~\cite{TODO toffola17} searches for literals inside of the PUT and derives suitable inputs.
PULSAR\cite{TODO wressnegger15} infers a network protocol from a set of captured network packages generated by the PUT.

Encoder-based models interpret the PUT as a decoder for inputs generated by an existing encoder.

\subsubsection{Mutation-based input generation}

Generating purely random inputs is not an effective strategy.
In general, a PUT contains checks that reject invalid inputs, making it
difficult to reach code paths that violate a given security policy~\cite{manes_art_2019}.
This motivates the use of a \emph{seed corpus}, a set of well-structured inputs
to the PUT like valid network packages or files.
The seeds inside of the seed corpus are then successively mutated with the goal
of triggering abnormal behavior.
Mutation operators are divided into four sets:

\begin{enumerate}
  \item \emph{Bit-flipping:}
    Bit flipping is governed by a \emph{mutation ratio}, the fraction of bits to
    flip per execution.

  \item \emph{Arithmetic mutation:}
    Arithmetic mutation reinterprets byte sequences as integers and perturbs them.

  \item \emph{Block-based mutation:}
    Block-based mutation inserts, deletes, replaces, permutes or resizes byte blocks,
    and splices blocks across seeds.

  \item \emph{Dictionary-based mutations:}
    Dictionary-based mutations substitute values with other values of high semantic
    weight. These usually include numbers like 0, or -1, or format strings for
    string substitution.
\end{enumerate}

\subsection{Coverage-Guided Greybox Fuzzing}\label{sec:fuzzing-cgf}

This work will focus primarily on coverage-guided greybox fuzzing, \emph{CFG} for short.
We will outline the general approach following Böhme et al.~\cite{TODO cite boehme17}.

\subsubsection{Overview}

% TODO add algorithm from boehme17, $2.1

A coverage-guided greybox fuzzer in general uses mutational input generation.
For this, it keeps a set of test cases $T$ as part of its runtime state.
This set is initialized from a seed corpus $S$.
The fuzzer then selects an input $t\in T$ and derives a new input $t^\prime$.
If the fuzzer deems the execution of the PUT under $t^\prime$ to be interesting, $t^\prime$
is then added to the test case set $T$.

The fuzzer must therefore solve mainly three problems: selecting an appropriate
test case, applying appropriate mutations and deciding whether to retain the
mutated input or not based on PUT behavior.

\subsubsection{Test case selection}

% TODO

\subsubsection{Test case mutation}

% TODO

\subsubsection{Test case retention}

% TODO: Describe different coverage metrics in detail and outline how this connects
% to instrumentation.

\subsubsection{AFL++}

One famous example of a coverage-guided greybox fuzzer is AFL++~\cite{fioraldi_afl_nodate}.
It is a mutation-based fuzzer that iteratively improves the seed corpus through
a series of mutations and selections.
Mutations include operations like bit flips, substitutions and splicing.
Newly generated inputs are then evaluated and added to the seed corpus for further
mutation and selection if they produce new coverage.
The success of a greybox fuzzer is strongly dependent on a high-quality initial
corpus and appropriate seed mutation and selection strategies.

% TODO add AFL algorithm (selection, mutation, etc.), maybe with reference to previous,
% more general fuzzing figure.

\begin{itemize}
  \item Frame CGF as the dominant modern design, popularized by AFL and refined by
    AFL++~\cite{fioraldi_afl_nodate} \todo{cite: original AFL (Zalewski)}.
  \item Core feedback loop, generalizing the excerpt above (describe as
    algorithm/figure)~\cite{manes_art_2019,fioraldi_afl_nodate}:
    \begin{itemize}
      \item pick a seed from the corpus (seed scheduling / energy assignment),
      \item mutate it into a batch of test cases,
      \item execute the instrumented target on each test case,
      \item bug oracle check: crash/hang $\rightarrow$ report,
      \item coverage check: new coverage $\rightarrow$ add test case to corpus,
      \item repeat.
    \end{itemize}
  \item Evolutionary interpretation: corpus = population, coverage novelty = fitness function,
    mutation/splicing = genetic operators~\cite{she_neuzz_2019,manes_art_2019}.
  \item Coverage signal design:
    \begin{itemize}
      \item granularities: basic-block vs.\ edge/branch coverage, with bucketed hit
        counts~\cite{fioraldi_afl_nodate},
      \item AFL-style fixed-size bitmap indexed by hashed edge IDs: fast but suffers hash
        collisions that alias distinct edges
        \todo{cite: CollAFL (Gan et al.\ 2018) for the collision analysis and fix},
      \item alternatives: hardware branch tracing (Intel PT) avoids compile-time
        instrumentation~\cite{zhang_ptfuzz_2018},
      \item relevant later: MOGI's basic-block hit tracker uses the same hash-into-buckets idea
        (\cref{chap:mogi-setup}).
    \end{itemize}
  \item Instrumentation options: compile-time insertion, static binary rewriting, or dynamic
    translation under emulation for closed-source binaries~\cite{fioraldi_afl_nodate}
    \todo{cite: QEMU (Bellard 2005) for dynamic binary translation}.
  \item Bug oracles beyond crashes: sanitizers (e.g.\ AddressSanitizer) surface silent memory
    corruption, at significant runtime overhead~\cite{zhang_debloating_nodate}
    \todo{cite: original AddressSanitizer paper (Serebryany et al.\ 2012)}.
  \item Engineering for throughput: fork-server and persistent modes, snapshotting; scaling across
    cores is itself nontrivial (program-adaptive parallel fuzzing)~\cite{zhou_kraken_2025}.
  \item Scheduling refinements: power schedules biasing energy toward rarely exercised paths
    \todo{cite: AFLFast (Böhme et al.\ 2016)}; learned seed scheduling~\cite{wang_reinforcement_2021}.
  \item Directed greybox fuzzing: optimize distance to chosen code sites instead of global
    coverage --- e.g.\ for patch testing or reaching suspected-vulnerable
    code~\cite{chen_critical_nodate,kwon_optimization-directed_2025}
    \todo{cite: AFLGo (Böhme et al.\ 2017) as the canonical directed greybox fuzzer}.
\end{itemize}

\subsection{Strengths and Limitations of Coverage-Guided Fuzzing}\label{sec:fuzzing-limitations}

\begin{itemize}
  \item Strengths (why CGF won):
    \begin{itemize}
      \item simple, robust, almost no per-target manual effort~\cite{manes_art_2019},
      \item scales linearly with compute; embarrassingly parallel in
        principle~\cite{zhou_kraken_2025},
      \item empirically the most productive bug-finding technique of the last
        decade~\cite{manes_art_2019,fioraldi_afl_nodate}.
    \end{itemize}
  \item Limitations --- each one motivates ``smarter'' input generation and frames a design goal of
    this thesis:
    \begin{itemize}
      \item \emph{Coverage plateaus}: random mutation practically never satisfies narrow
        constraints (magic bytes, checksums, length fields), so campaigns stall on frontier
        branches~\cite{stephens_driller_2016,she_neuzz_2019}.
      \item \emph{Coverage is only a proxy}: maximizing edges is not the same as triggering bugs;
        state- and data-dependent bugs need specific values in specific contexts, not just new
        edges~\cite{chen_critical_nodate}.
      \item \emph{Blind mutation wastes the execution budget}: most mutations are irrelevant to
        control flow; which bytes and operators matter is input- and target-dependent
        ~\cite{lyu_mopt_nodate,binosi_rainfuzz_2023}.
      \item \emph{Feedback aliasing}: bitmap hash collisions and coarse hit-count bucketing hide
        genuinely new behavior \todo{cite: CollAFL}.
      \item \emph{Statelessness}: the classic loop assumes one self-contained input per fresh
        execution --- breaks down for stateful/reactive targets (\cref{sec:fuzzing-stateful}).
      \item \emph{Evaluation is tricky}: campaigns are highly stochastic; sound comparisons need
        many trials and statistical testing \todo{cite: Klees et al.\ 2018, ``Evaluating Fuzz
        Testing''} (also needed for \cref{sec:analysis-evaluation}).
    \end{itemize}
  \item Research responses (survey briefly, as context for the RL approach):
    \begin{itemize}
      \item data-flow/taint guidance to identify influential input bytes
        (GreyOne)~\cite{gan_greyone_nodate},
      \item neural surrogates: smooth the coverage function and follow gradients (NEUZZ,
        MTFuzz)~\cite{she_neuzz_2019,she_mtfuzz_2020}, deep-learning path prediction
        (NeuFuzz)~\cite{wang_neufuzz_2019},
      \item hybrid concolic execution (Driller)~\cite{stephens_driller_2016},
      \item machine learning across all loop stages --- seed selection, mutation, scheduling ---
        surveyed in~\cite{wang_systematic_2020,qiu_survey_2025}.
    \end{itemize}
\end{itemize}

\subsection{Stateful and Protocol Fuzzing}\label{sec:fuzzing-stateful}

\begin{itemize}
  \item Many high-value targets are reactive servers speaking a protocol (network services,
    industrial control systems, embedded devices): behavior depends on \emph{session state}
    accumulated over a message exchange~\cite{ang_quic-fuzz_2025,borcherding_bandits_2023}.
  \item Consequences for the classic single-input model:
    \begin{itemize}
      \item a test case is a \emph{sequence} of messages, not one blob; ordering and protocol
        context matter,
      \item deep states are reachable only through valid message prefixes; mutating a late message
        is useless if an early one breaks the session,
      \item the same code can behave differently in different states, so pure edge coverage
        under-represents state coverage~\cite{borcherding_bandits_2023},
      \item state is invisible: the fuzzer must infer or track it (response codes, annotated state
        variables) \todo{cite: AFLNet (Pham et al.\ 2020) and SGFuzz (Ba et al.\ 2022) --- the
        standard stateful greybox fuzzers; both are baselines the thesis should discuss}.
      \item state-awareness matters in adjacent domains too, e.g.\ black-box web
        scanning~\cite{doupe_enemy_2012}.
    \end{itemize}
  \item Extra decision problem on top of the CGF loop: \emph{which state to fuzz next} ---
    an exploration/exploitation trade-off, already modeled as a multi-armed bandit in prior
    work~\cite{borcherding_bandits_2023}; direct bridge to \cref{sec:reinforcement-learning}.
  \item Practical complications for encrypted/checksummed transports (e.g.\ QUIC): the fuzzer must
    sit inside the crypto boundary to mutate meaningfully~\cite{ang_quic-fuzz_2025}.
  \item Position of this thesis: protocol-level actions against a stateful target are exactly the
    action space handed to the RL agent (\cref{chap:efficientzero-fuzzing}).
\end{itemize}

\subsection{Emulator-Based Fuzzing}\label{sec:fuzzing-emulation}

\begin{itemize}
  \item Motivation from embedded/firmware targets: fuzzing on the device is slow and, worse,
    faults are often \emph{silent} --- without MMU protection or sanitizers, memory corruption
    frequently does not crash (``what you corrupt is not what you
    crash'')~\cite{muench_what_2018}.
  \item Rehosting the target in an emulator restores the missing capabilities~\cite{muench_what_2018}
    \todo{cite: QEMU (Bellard 2005); Unicorn engine; optionally a rehosting survey (e.g.\ Wright
    et al.\ 2021 or Fasano et al.\ 2021)}:
    \begin{itemize}
      \item full observability: coverage and memory-access instrumentation without source access,
      \item determinism and replayability of executions,
      \item fine-grained fault detection via emulator-level checks,
      \item \emph{snapshots}: capture full machine state once, restore cheaply before every test.
    \end{itemize}
  \item Snapshot restore gives an exact, cheap \emph{reset} of the target to a known state ---
    for stateful targets this is what makes per-episode experimentation tractable, and it is
    precisely the environment-reset primitive that reinforcement learning assumes
    (\cref{sec:reinforcement-learning}); realized here by the MOGI framework
    (\cref{chap:mogi-setup}).
  \item Cost: emulation overhead reduces raw executions per second versus native execution ---
    sharpens the need to make every execution count (sample efficiency,
    \cref{chap:efficientzero}).
\end{itemize}
